% ===================================================================
%  TABEL Nr. 15 — Die mark. — Die eetware.
%  Meme regime que les precedents. Tableau a UNE SEULE COLONNE.
%
%  UN RENVOI DE CE TABLEAU EST CORRIGE SUR LA PAGE DE LECTURE, DANS
%  LA COLONNE FRANCAISE SEULE : elle imprime (32) au 13 la ou l'ido
%  lit (82), les petits pains. La correction vit dans
%  gravuri/korekti.json et ne touche aucune source.
%
%  LE TABLEAU EST UN INVENTAIRE, et c'est la sa difficulte : cent
%  cinq renvois, presque tous des denrees, et beaucoup se ressemblent
%  assez pour tomber sur le meme mot si on n'y prend pas garde.
%  « pruni » (53) et « sika pruni » (65) sont deux renvois ;
%  « korbi » (70) et « provizo-reti » (73) aussi ; et le tableau
%  distingue « bucho-servisto » (75), qui pousse la brouette, de
%  « bucho-garsono » (93), qui decoupe — slagtersbediende et
%  slagterskneg.
%
%  ET LES TROIS VIANDES DOIVENT RESTER TROIS : « muton-karno » (90),
%  « bov-karno » (94) et « bovyun-karno » (92) — skaapvleis,
%  beesvleis, kalfsvleis.
% ===================================================================

%%K t15-apar-1 apar
\VUsaut{4.0mm}
\VUtitre{15.0pt}{60}{TABEL Nr. 15}
\VUsaut{3.0mm}

%%K t15-tit-1 sub
\VUpk{11.4pt}{40}{Die mark. --- Die eetware.}
\VUsaut{3.0mm}

%%K t15-01-1 p
1. --- Die meeste handelaars wat vir ons die ware lewer wat vir ons voeding nodig is,
versamel hulle onder die \VUgras{hal} \textsuperscript{(1)} van die \VUgras{oordekte
mark} of in die aangrensende strate.

%%K t15-02-1 p
2. --- Die \VUgras{varkslagter} \textsuperscript{(2)}, wat \VUgras{ham}
\textsuperscript{(3)}, \VUgras{bloedwors} \textsuperscript{(4)}, \VUgras{worsie}
\textsuperscript{(5)}, \VUgras{pensworsie} \textsuperscript{(6)} en \VUgras{spek}
\textsuperscript{(7)} verkoop, staan naby die \VUgras{botter- en kaasverkoopster}
\textsuperscript{(8)}, wie se uitstalling voorsien is van \VUgras{Hollandse kase}
\textsuperscript{(9)} met 'n rooi kors, van \VUgras{Camembertkase}
\textsuperscript{(10)}, van \VUgras{Gruyèreskywe} \textsuperscript{(11)} en van vars
\VUgras{botterblokke} \textsuperscript{(12)}.

%%K t15-03-1 p
3. --- Voor haar bied 'n \VUgras{groenteverkoopster} \textsuperscript{(13)} aan haar
klante mooi \VUgras{tamaties} \textsuperscript{(14)}, \VUgras{blomkole}
\textsuperscript{(15)}, \VUgras{slaai} \textsuperscript{(16)}, \VUgras{kopkole}
\textsuperscript{(17)}, \VUgras{murgpampoentjies} \textsuperscript{(19)},
\VUgras{spanspekke} \textsuperscript{(18)} en selfs \VUgras{sampioene}
\textsuperscript{(20)} aan. Maar 'n \VUgras{bierkneg}, wat sy \VUgras{afleweringswa}
\textsuperscript{(21)} gelaai met \VUgras{biervate} \textsuperscript{(22)} juis voor haar
uitstalling laat stilhou het, belet haar klante om te nader. 'n Tuinierster bring vir
haar 'n groot mandjie \VUgras{artisjokke} \textsuperscript{(23)}.

%%K t15-tit-2 sub
\VUsaut{4.0mm}
\VUpk{10.2pt}{40}{Verkoop en koop.}
\VUsaut{3.0mm}

%%K t15-04-1 p
4. --- Op die mark is die bedrywigheid groot, want die uur van die \VUgras{veiling}
\textsuperscript{(24)} het aangebreek. Die \VUgras{huisvroue} \textsuperscript{(26)},
wat daarheen vir hulle inkopies kom, hou in die verbygaan stil voor die stalletjie van
die \VUgras{pensverkoopster} \textsuperscript{(25)}, om \VUgras{pens} of 'n
\VUgras{kalfskop} te koop.

%%K t15-05-1 p
5. --- 'n Vet \VUgras{kok} \textsuperscript{(27)} ding af op 'n baie mooi \VUgras{tuna}
by die \VUgras{visverkoopster} \textsuperscript{(28)} wat na haar goeddunke haar
(verkoop)pryse verhoog. Sy het 'n groot hoeveelheid \VUgras{palings, tonge, forelle} en
\VUgras{karpe} gekry. Haar \VUgras{kabeljou} en \VUgras{mossels} is baie vars.

%%K t15-tit-3 sub
\VUsaut{4.0mm}
\VUpk{10.2pt}{40}{Goedkoop en duur ware.}
\VUsaut{3.0mm}

%%K t15-06-1 p
6. --- Die \VUgras{pluimveeverkoopster} \textsuperscript{(29)}, wat naby haar gevestig
is, is baie nougeset. Wanneer sy 'n \VUgras{kalkoen, 'n gans, 'n eend} of 'n gewone
\VUgras{hoender} verkoop, vra sy net die regte prys.

%%K t15-07-1 p
7. --- Op die hoek van die plein, op die eerste verdieping, is sedert kort 'n
\VUgras{linneverkoper} \textsuperscript{(30)} gevestig by wie die koopsters altyd seker
is om 'n baie mooi verskeidenheid \VUgras{tafeldoeke,} \VUgras{servette, voorskote} en
\VUgras{handdoeke} te vind. Op dieselfde verdieping het 'n \VUgras{messemaker}
\textsuperscript{(31)} sy werkplek gevestig. Hy slyp die \VUgras{messe} van die
verkoopsters in die hal en maak vir die tuiniers \VUgras{snoeimesse} van die hardste
\VUgras{staal.}

%%K t15-tit-4 sub
\VUsaut{4.0mm}
\VUpk{10.2pt}{40}{Die kruidenierswinkel.}
\VUsaut{3.0mm}

%%K t15-08-1 p
8. --- Onder sy werkplek is sedert kort 'n groot lewensmiddelewinkel oopgemaak. Dit is
al nie meer die eenvoudige en onbelangrike \VUgras{kruidenierswinkel}
\textsuperscript{(32)} van vroeër nie, wat net die gangbare ware verkoop het,
\VUgras{tee} \textsuperscript{(33)}, \VUgras{sjokolade} \textsuperscript{(34)},
\VUgras{broodsuiker} \textsuperscript{(37)} en \VUgras{seep} \textsuperscript{(38)}. Daar
verkoop 'n mens enige ware, \VUgras{bros koekies}, \VUgras{konserwe}
\textsuperscript{(35)}, \VUgras{likeure} \textsuperscript{(36)}, \VUgras{rum,}
\VUgras{konjak, kirsch, anysdrank} en \VUgras{curaçao.} Daar vind 'n mens wild:
\VUgras{haas} \textsuperscript{(39)}, \VUgras{lysters} \textsuperscript{(40)},
\VUgras{patryse} \textsuperscript{(41)}, \VUgras{kwartels} \textsuperscript{(42)},
\VUgras{wilde eend} \textsuperscript{(43)} of \VUgras{fisant} \textsuperscript{(44)}, net
so maklik soos uitheemse vrugte soos \VUgras{pisangs} \textsuperscript{(46)} en
\VUgras{pynappels.}

%%K t15-09-1 p
9. --- 'n Baie hulpvaardige \VUgras{kruidenierskneg} \textsuperscript{(45)} is gereed om
te gee wat 'n mens verlang: \VUgras{konfyt} \textsuperscript{(47)} van aalbessies of van
kwepers, \VUgras{vet} \textsuperscript{(48)}, \VUgras{agurkies} \textsuperscript{(49)} of
gesoute \VUgras{sardientjies} \textsuperscript{(50)}.

%%K t15-09-2 p
Dit is baie gerieflik vir die \VUgras{kliënte} \textsuperscript{(51)}, wat naas die
gangbaarste ware die seldsaamste eerstelinge en die smaaklikste vrugte vind: sappige
\VUgras{pere} \textsuperscript{(52)}, \VUgras{pruime} \textsuperscript{(53)},
\VUgras{druiwe} \textsuperscript{(54)}, \VUgras{perskes} \textsuperscript{(55)},
\VUgras{amandels} \textsuperscript{(56)} en \VUgras{neute} \textsuperscript{(57)}.

%%K t15-tit-5 sub
\VUsaut{4.0mm}
\VUpk{10.2pt}{40}{Die klantekring.}
\VUsaut{3.0mm}

%%K t15-10-1 p
10. --- Die \VUgras{rys} \textsuperscript{(58)} en die \VUgras{lensies}
\textsuperscript{(59)} staan langs die kis met gerookte \VUgras{harings}
\textsuperscript{(60)}; die \VUgras{aartappel} \textsuperscript{(61)} en die
\VUgras{boontjie} \textsuperscript{(63)} is naby die \VUgras{appels}
\textsuperscript{(62)}, \VUgras{vye} \textsuperscript{(64)}, \VUgras{gedroogde pruime}
\textsuperscript{(65)} en \VUgras{haselneute} \textsuperscript{(66)}. 'n Mens vind in
daardie winkel vars \VUgras{eiers} \textsuperscript{(67)} en \VUgras{olywe}
\textsuperscript{(68)}; daarom word \VUgras{mandjies} \textsuperscript{(70)} en
\VUgras{boodskapnette} \textsuperscript{(73)} blitsvinnig gevul.

%%K t15-11-1 p
11. --- Die \VUgras{koffie} \textsuperscript{(72)} van daardie uitstekende firma het
onder die \VUgras{diensmeisies} \textsuperscript{(69)} en die \VUgras{kokkinne} 'n
verdiende naam verwerf, want die ou \VUgras{kruidenier} \textsuperscript{(71)} rooster
dit self met die grootste sorg.

%%K t15-tit-6 sub
\VUsaut{4.0mm}
\VUpk{10.2pt}{40}{Die skouspele.}
\VUsaut{3.0mm}

%%K t15-12-1 p
12. --- Nie almal gaan na die mark om hulle te bevoorraad nie. In die skilderagtige
verkeer op die straatjie wat na die hal lei, sien 'n mens die mees uiteenlopende mense.
Naas 'n vriendelike \VUgras{slagtersbediende} \textsuperscript{(75)}, wat 'n
\VUgras{kruiwa} \textsuperscript{(74)} voor hom uitstoot, is daar twee soldate met verlof
wat heeltemal trots is om hulle blinkende uniform te wys: die \VUgras{husaar}
\textsuperscript{(76)} met 'n blou \VUgras{dolman} en 'n \VUgras{sjako met pluim}, en die
\VUgras{soeaafkorporaal,} met wye broek en die kop bedek met 'n \VUgras{fes.}

%%K t15-13-1 p
13. --- Die \VUgras{bakker} \textsuperscript{(80)}, wat op die drumpel van die
\VUgras{bakkery} \textsuperscript{(78)} staan, het pas die deeg van sy \VUgras{brood}
\textsuperscript{(79)} geknie en uit die oond die \VUgras{croissants}
\textsuperscript{(81)}, die \VUgras{broodjies} \textsuperscript{(82)} en die
\VUgras{ronde brode} \textsuperscript{(83)} gehaal wat in die uitstalvenster is.

%%K t15-14-1 p
14. --- Bo die bakkery woon 'n bekwame \VUgras{linnenaaister} \textsuperscript{(84)} wat
etlike \VUgras{borduursters} \textsuperscript{(85)} in diens het. Naas haar woon 'n
\VUgras{wasvrou-strykster} \textsuperscript{(86)}, wat die \VUgras{linnegoed}
\textsuperscript{(88)} styf nadat sy dit gewas en gedroog het, en wat dit met 'n
\VUgras{strykyster} \textsuperscript{(87)} glad maak.

%%K t15-tit-7 sub
\VUsaut{4.0mm}
\VUpk{10.2pt}{40}{Vleis, vis en groente.}
\VUsaut{3.0mm}

%%K t15-15-1 p
15. --- In die \VUgras{slaghuis} \textsuperscript{(89)}, wat op die grondverdieping lê,
verkoop 'n mens allerhande \VUgras{vleis, skaapvleis} \textsuperscript{(90)},
\VUgras{beesvleis} \textsuperscript{(94)} of \VUgras{kalfsvleis} \textsuperscript{(92)}
as \VUgras{skaapboude, biefstukke, gebraad} en \VUgras{kotelette.} Die
\VUgras{slagterskneg} \textsuperscript{(93)} sny al daardie vleis op en sit die
\VUgras{murgbene} \textsuperscript{(96)} apart. By die deur van die slaghuis is 'n
\VUgras{oesterverkoopster} \textsuperscript{(97)} wat \VUgras{oesters}
\textsuperscript{(98)} verkoop. Sy maak hulle oop as 'n mens dit verlang.

%%K t15-16-1 p
16. --- Al daardie handelaars betaal 'n patent of 'n staangeld; daarom vervolg die
\VUgras{polisieman} \textsuperscript{(100)} sonder genade die arme
\VUgras{groentesmouse} \textsuperscript{(99)} wat met hulle meeding deur met hulle
karretjies \VUgras{wortels, radysies, rape, prei} of \VUgras{aspersiebondels, vars
ertjies, spinasie, suring, bronkors} en \VUgras{pietersielie} rond te vent.

%%K t15-tit-8 sub
\VUsaut{4.0mm}
\VUpk{10.2pt}{40}{Pas op vir die polisieman!}
\VUsaut{3.0mm}

%%K t15-17-1 p
17. --- Die seun wat \VUgras{knoffel} \textsuperscript{(101)} en \VUgras{uie} verkoop,
weet baie goed dat ook hy sy ware nie naby die mark mag verkoop nie. Hy vlug en loop
daarby gevaar om die mandjie \VUgras{krappe,} \VUgras{varswaterkrewe} en
\VUgras{garnale} van die \VUgras{verkoper} \textsuperscript{(102)} van
\VUgras{langoeste} en \VUgras{seekrewe} om te gooi.

%%K t15-18-1 p
18. --- Almal beweeg hulle en maak baie lawaai; maar dit belet nie om die stem van die
rondtrekkende \VUgras{glasmaker} \textsuperscript{(103)} duidelik te hoor nie, of die
kreet van die \VUgras{kolehandelaar} \textsuperscript{(104)} wat pas 'n paar oomblikke sy
karretjie verlaat het om 'n \VUgras{sak steenkool} \textsuperscript{(105)} te lewer.
\VUsaut{6.0mm}
\VUfilet{22mm}
